\chapter{Análisis de Impacto:}
%TODO: ANÁLISIS DE IMPACTO.

Con este desarrollo software se espera mejorar distintas áreas del laboratorio, principalmente las siguientes. La primera es la gestión del inventario de dispositivos, máquinas y servicios, lo que permitirá al administrador disponer de una visión actualizada en tiempo real. La segunda es la trazabilidad a través del sistema de peticiones, mediante el cual el jefe de laboratorio deberá aprobar todos los cambios solicitados, manteniéndose así informado en todo momento, así como los supervisores de cada usuario. Por último, en el ámbito de la gestión de usuarios, se establece una única fuente de verdad (single source of truth) que el administrador mantiene actualizada en tiempo real.

Asimismo, el personal investigador dispondrá de un entorno privado, seguro y virtualizado con acceso a los distintos recursos de la máquina, facilitando así la realización de sus tareas. Dicho entorno incorpora todas las funcionalidades de Visual Studio Code, incluida la posibilidad de instalar extensiones, gracias a la infraestructura de code-server más las mejoras incorporadas.
